% Generated by scripts/render_mcr_editions.py; edit the shared register.
\paragraph{Operating constraints and evidence.} 
% Shared source: ANT-THM-FR-001
The thermal subsystem must keep each selected component within traceable operating and survival temperature limits for nominal and Safe Mode cases; numerical limits/node mapping remain open allocations. One history: internal −14.061 to21.159°C; no component compliance limits.
% Shared source: ANT-THM-FR-002
The thermal assessment must bound hot, cold and eclipse temperatures using mission-compatible attitude and mode-dependent internal heat loads, including startup, science, Relay and safe recovery. Supplied extrema are one history; 1.9–4.9 W schedule is independent of DITL.

% Shared source: ANT-THM-FR-003
The thermal interfaces must allocate battery charge/discharge temperature constraints and payload temperature-dependent calibration/stability limits. Both numerical allocations remain open allocations. No selected-part limits or payload thermal response characterization.
% Shared source: ANT-THM-FR-004
The thermal-mechanical design must keep gradients and resulting distortions/stresses within allocated alignment and material limits; numerical limits remain open allocations. 31.704°C face spread is a node difference, not structural stress.

% Shared source: ANT-STR-FR-001
The structure/integration budget must fit the selected 3 U deployer envelope and keep the complete design mass, including all allocations and contingency, at or below 3.0 kg. 1536.3 g is incomplete BOM; provider ICD, payload and deployment allocations open.
% Shared source: ANT-STR-FR-002
The structure must withstand applicable launch loads and meet provider modal, vibration and shock requirements without loss of function or interface integrity. Numerical load/frequency limits remain open allocations. Launch-provider load envelope and structural/modal analysis are not available.

% Shared source: ANT-STR-FR-003
The solar deployment/retention mechanism must retain both panels safely through launch and release/verify deployment without collision, harness damage or loss of the required generation geometry. Deployment architecture is defined; mechanism loads, clearances and retention life are not assessed.
% Shared source: ANT-STR-FR-004
The mechanical integration must preserve FPDU optical/readout interfaces and body-frame clearances through assembly, thermal cycling and operational life. Active material dimensions exclude coating, gaps, sensors and housing.

\paragraph{Feasibility assessment.} The thermal model predicts a finite temperature envelope for its evaluated heat schedule, providing a basis for component selection and thermal design. Component compliance and structural feasibility are not yet assessed against selected-part limits, a complete mass allocation and launch-provider loads.
